\section{Supplement on quasi-coherent sheaves}
\label{section:I.9}

\subsection{Tensor product of quasi-coherent sheaves.}
\label{subsection:I.9.1}

\begin{proposition}[9.1.1]
\label{I.9.1.1}
Let $X$ be a prescheme (resp. a locally Noetherian prescheme). Let
$\mathscr{F}$ and $\mathscr{G}$ be two quasi-coherent (resp. coherent)
$\mathscr{O}_X$-modules; then
$\mathscr{F}\otimes_{\mathscr{O}_X}\mathscr{G}$ is quasi-coherent (resp.
coherent), and of finite type if $\mathscr{F}$ and $\mathscr{G}$ are of
finite type. If $\mathscr{F}$ admits a finite presentation and if
$\mathscr{G}$ is quasi-coherent (resp. coherent), then
$\mathscr{H}\!om_{\mathscr{O}_X}(\mathscr{F},\mathscr{G})$ is quasi-coherent
(resp. coherent).
\end{proposition}

The question being local, we may suppose $X$ affine (resp. affine
Noetherian); moreover, if $\mathscr{F}$ is coherent, we may suppose that it
is the cokernel of a homomorphism
$\mathscr{O}_X^m\to\mathscr{O}_X^n$. The assertions concerning quasi-coherent
sheaves then follow from
(\hyperref[I.1.3.12]{1.3.12}) and
(\hyperref[I.1.3.9]{1.3.9}); the assertions concerning coherent
sheaves follow from (\hyperref[I.1.5.1]{1.5.1}) and from the fact that, if
$M$ and $N$ are modules of finite type over a Noetherian ring $A$,
$M\otimes_A N$ and $\operatorname{Hom}_A(M,N)$ are $A$-modules of finite
type.

\begin{definition}[9.1.2]
\label{I.9.1.2}
Let $X$, $Y$ be two $S$-preschemes, $p$, $q$ the projections of
$X\times_S Y$, and $\mathscr{F}$ (resp. $\mathscr{G}$) a quasi-coherent
$\mathscr{O}_X$-module (resp. $\mathscr{O}_Y$-module). We call tensor
product of $\mathscr{F}$ and $\mathscr{G}$ over $\mathscr{O}_S$ (or over
$S$), and denote by
$\mathscr{F}\otimes_{\mathscr{O}_S}\mathscr{G}$ (or
$\mathscr{F}\otimes_S\mathscr{G}$), the tensor product
$p^*(\mathscr{F})\otimes_{\mathscr{O}_{X\times_S Y}}q^*(\mathscr{G})$ on the
prescheme $X\times_S Y$.
\end{definition}

If $X_i$ ($1\leq i\leq n$) are $S$-preschemes and $\mathscr{F}_i$ is a
quasi-coherent $\mathscr{O}_{X_i}$-module ($1\leq i\leq n$), we define in
the same way the tensor product
$\mathscr{F}_1\otimes_S\mathscr{F}_2\otimes_S\cdots\otimes_S\mathscr{F}_n$
on the prescheme
$Z=X_1\times_S X_2\times_S\cdots\times_S X_n$; it is a
\emph{quasi-coherent} $\mathscr{O}_Z$-module by virtue of
(\hyperref[I.9.1.1]{9.1.1}) and
(\hyperref[0.5.1.4]{0, 5.1.4}); it is \emph{coherent} if the
$\mathscr{F}_i$ are and if $Z$ is \emph{locally Noetherian}, by virtue
of (\hyperref[I.9.1.1]{9.1.1}),
(\hyperref[0.5.3.11]{0, 5.3.11}) and
(\hyperref[I.6.1.1]{6.1.1}).

Note that if we take $X=Y=S$, definition
(\hyperref[I.9.1.2]{9.1.2}) gives back the tensor product of
$\mathscr{O}_S$-modules. Moreover, since
$q^*(\mathscr{O}_Y)=\mathscr{O}_{X\times_S Y}$
(\hyperref[0.4.3.4]{0, 4.3.4}), the product
$\mathscr{F}\otimes_S\mathscr{O}_Y$ is canonically identified with
$p^*(\mathscr{F})$, and likewise
$\mathscr{O}_X\otimes_S\mathscr{G}$ is canonically identified with
$q^*(\mathscr{G})$. More particularly, if we take $Y=S$ and denote by $f$
the structural morphism $X\to Y$, we have
$\mathscr{O}_X\otimes_Y\mathscr{G}=f^*(\mathscr{G})$: the ordinary tensor
product and the inverse image thus appear as particular cases of the
general tensor product.

Definition (\hyperref[I.9.1.2]{9.1.2}) immediately implies that, for
fixed $X$ and $Y$, $\mathscr{F}\otimes_S\mathscr{G}$ is a
\emph{covariant, additive and right exact bifunctor} in $\mathscr{F}$ and
$\mathscr{G}$.

\begin{proposition}[9.1.3]
\label{I.9.1.3}
Let $S$, $X$, $Y$ be three affine schemes with respective rings $A$, $B$,
$C$, $B$ and $C$ being $A$-algebras. Let $M$ (resp. $N$) be a $B$-module
(resp. $C$-module), and $\mathscr{F}=\widetilde{M}$ (resp.
$\mathscr{G}=\widetilde{N}$) the associated quasi-coherent sheaf; then
$\mathscr{F}\otimes_S\mathscr{G}$ is canonically isomorphic to the sheaf
associated to the $(B\otimes_A C)$-module $M\otimes_A N$.
\end{proposition}

\begin{proof}
\oldpage[I]{170}
Indeed, by virtue of (\hyperref[I.1.6.5]{1.6.5}),
$\mathscr{F}\otimes_S\mathscr{G}$ is canonically isomorphic to the sheaf
associated to the $(B\otimes_A C)$-module
\[
  \bigl(M\otimes_B(B\otimes_A C)\bigr)
  \otimes_{B\otimes_A C}
  \bigl((B\otimes_A C)\otimes_C N\bigr)
\]
and, owing to the canonical isomorphisms between tensor products, the
latter is isomorphic to
\[
  M\otimes_B(B\otimes_A C)\otimes_C N
  =(M\otimes_B B)\otimes_A(C\otimes_C N)=M\otimes_A N.
\]
\end{proof}

\begin{proposition}[9.1.4]
\label{I.9.1.4}
Let $f:T\to X$, $g:T\to Y$ be two $S$-morphisms, and $\mathscr{F}$ (resp.
$\mathscr{G}$) a quasi-coherent $\mathscr{O}_X$-module (resp.
$\mathscr{O}_Y$-module). Then we have
\[
  (f,g)_S^*(\mathscr{F}\otimes_S\mathscr{G})
  =f^*(\mathscr{F})\otimes_{\mathscr{O}_T}g^*(\mathscr{G}).
\]
\end{proposition}

\begin{proof}
If $p$, $q$ are the projections of $X\times_S Y$, the formula indeed
follows from the relations $(f,g)_S^*\circ p^*=f^*$ and
$(f,g)_S^*\circ q^*=g^*$ (\hyperref[0.3.5.5]{0, 3.5.5}), and from the
fact that an inverse image of a tensor product of algebraic sheaves is
the tensor product of their inverse images
(\hyperref[0.4.3.3]{0, 4.3.3}).
\end{proof}

\begin{corollary}[9.1.5]
\label{I.9.1.5}
Let $f:X\to X'$, $g:Y\to Y'$ be two $S$-morphisms, and $\mathscr{F}'$
(resp. $\mathscr{G}'$) a quasi-coherent $\mathscr{O}_{X'}$-module (resp.
$\mathscr{O}_{Y'}$-module). Then we have
\[
  (f\times_S g)^*(\mathscr{F}'\otimes_S\mathscr{G}')
  =f^*(\mathscr{F}')\otimes_S g^*(\mathscr{G}').
\]
\end{corollary}

\begin{proof}
This follows from (\hyperref[I.9.1.4]{9.1.4}) and from the fact that
$f\times_S g=(f\circ p,g\circ q)_S$, $p$ and $q$ being the projections of
$X\times_S Y$.
\end{proof}

\begin{corollary}[9.1.6]
\label{I.9.1.6}
Let $X$, $Y$, $Z$ be three $S$-preschemes, and $\mathscr{F}$ (resp.
$\mathscr{G}$, $\mathscr{H}$) a quasi-coherent $\mathscr{O}_X$-module
(resp. $\mathscr{O}_Y$-module, $\mathscr{O}_Z$-module); the sheaf
$\mathscr{F}\otimes_S\mathscr{G}\otimes_S\mathscr{H}$ is the inverse image
of $(\mathscr{F}\otimes_S\mathscr{G})\otimes_S\mathscr{H}$ under the
canonical isomorphism of $X\times_S Y\times_S Z$ onto
$(X\times_S Y)\times_S Z$.
\end{corollary}

\begin{proof}
Indeed, this isomorphism can be written $(p_1,p_2)_S\times_S p_3$,
denoting by $p_1$, $p_2$, $p_3$ the projections of $X\times_S Y\times_S Z$.

Likewise, the inverse image of $\mathscr{G}\otimes_S\mathscr{F}$ under the
canonical isomorphism of $X\times_S Y$ onto $Y\times_S X$ is
$\mathscr{F}\otimes_S\mathscr{G}$.
\end{proof}

\begin{corollary}[9.1.7]
\label{I.9.1.7}
If $X$ is an $S$-prescheme, every quasi-coherent $\mathscr{O}_X$-module
$\mathscr{F}$ is the inverse image of
$\mathscr{F}\otimes_S\mathscr{O}_S$ under the canonical isomorphism of $X$
onto $X\times_S S$ (\hyperref[I.3.3.3]{3.3.3}).
\end{corollary}

\begin{proof}
Indeed, this isomorphism is $(1_X,\varphi)_S$, where $\varphi$ is the
structural morphism $X\to S$, and the corollary follows from
(\hyperref[I.9.1.4]{9.1.4}) and from the fact that
$\varphi^*(\mathscr{O}_S)=\mathscr{O}_X$.
\end{proof}

\begin{env}[9.1.8]
\label{I.9.1.8}
Let $X$ be an $S$-prescheme, $\mathscr{F}$ a quasi-coherent
$\mathscr{O}_X$-module, and $\varphi:S'\to S$ a morphism; we denote by
$\mathscr{F}_{(\varphi)}$ or $\mathscr{F}_{(S')}$ the quasi-coherent
sheaf $\mathscr{F}\otimes_S\mathscr{O}_{S'}$ on
$X\times_S S'=X_{(\varphi)}=X_{(S')}$; thus
$\mathscr{F}_{(S')}=p^*(\mathscr{F})$, where $p$ is the projection
$X_{(S')}\to X$.
\end{env}

\begin{proposition}[9.1.9]
\label{I.9.1.9}
Let $\varphi':S''\to S'$ be a morphism. For every quasi-coherent
$\mathscr{O}_X$-module $\mathscr{F}$ on the $S$-prescheme $X$,
$(\mathscr{F}_{(\varphi)})_{(\varphi')}$ is the inverse image of
$\mathscr{F}_{(\varphi\circ\varphi')}$ under the canonical isomorphism
$(X_{(\varphi)})_{(\varphi')}\xrightarrow{\sim}
X_{(\varphi\circ\varphi')}$ (\hyperref[I.3.3.9]{3.3.9}).
\end{proposition}

\begin{proof}
This follows immediately from the definitions and from
(\hyperref[I.3.3.9]{3.3.9}), and can also be written
\[
  (\mathscr{F}\otimes_S\mathscr{O}_{S'})
  \otimes_{S'}\mathscr{O}_{S''}
  =\mathscr{F}\otimes_S\mathscr{O}_{S''}.
  \tag{9.1.9.1}\label{I.9.1.9.1}
\]
\end{proof}

\begin{proposition}[9.1.10]
\label{I.9.1.10}
Let $Y$ be an $S$-prescheme and $f:X\to Y$ an $S$-morphism. For every
quasi-coherent $\mathscr{O}_Y$-module $\mathscr{G}$ and every morphism
$S'\to S$, we have
\[
  (f_{(S')})^*(\mathscr{G}_{(S')})
  =(f^*(\mathscr{G}))_{(S')}.
\]
\end{proposition}

\begin{proof}
This follows immediately from the commutativity of the diagram
\oldpage[I]{171}
\[
  \xymatrix{
    X_{(S')}\ar[r]^{f_{(S')}}\ar[d] &
    Y_{(S')}\ar[d]\\
    X\ar[r]_f &
    Y
  }
\]
\end{proof}

\begin{corollary}[9.1.11]
\label{I.9.1.11}
Let $X$, $Y$ be two $S$-preschemes, and $\mathscr{F}$ (resp. $\mathscr{G}$)
a quasi-coherent $\mathscr{O}_X$-module (resp. $\mathscr{O}_Y$-module).
The inverse image of the sheaf
$(\mathscr{F}_{(S')}\otimes_{S'}\mathscr{G}_{(S')})$ under the canonical
isomorphism
$(X\times_S Y)_{(S')}\xrightarrow{\sim}
(X_{(S')})\times_{S'}(Y_{(S')})$
(\hyperref[I.3.3.10]{3.3.10}) is equal to
$(\mathscr{F}\otimes_S\mathscr{G})_{(S')}$.
\end{corollary}

\begin{proof}
If $p$, $q$ are the projections of $X\times_S Y$, the isomorphism in
question is none other than $(p_{(S')},q_{(S')})_{S'}$; the corollary
follows from Propositions (\hyperref[I.9.1.4]{9.1.4}) and
(\hyperref[I.9.1.10]{9.1.10}).
\end{proof}

\begin{proposition}[9.1.12]
\label{I.9.1.12}
With the notation of (\hyperref[I.9.1.2]{9.1.2}), let $z$ be a point of
$X\times_S Y$, $x=p(z)$, $y=q(z)$; the stalk
$(\mathscr{F}\otimes_S\mathscr{G})_z$ is isomorphic to
$(\mathscr{F}_x\otimes_{\mathscr{O}_x}\mathscr{O}_z)
\otimes_{\mathscr{O}_z}
(\mathscr{G}_y\otimes_{\mathscr{O}_y}\mathscr{O}_z)
=\mathscr{F}_x\otimes_{\mathscr{O}_x}\mathscr{O}_z
\otimes_{\mathscr{O}_y}\mathscr{G}_y$.
\end{proposition}

\begin{proof}
Since we can reduce to the affine case, the proposition follows from
formula (\hyperref[I.1.6.5.1]{1.6.5.1}).
\end{proof}

\begin{corollary}[9.1.13]
\label{I.9.1.13}
If $\mathscr{F}$ and $\mathscr{G}$ are of finite type, we have
\[
  \operatorname{Supp}(\mathscr{F}\otimes_S\mathscr{G})
  =p^{-1}(\operatorname{Supp}(\mathscr{F}))
  \cap q^{-1}(\operatorname{Supp}(\mathscr{G})).
\]
\end{corollary}

\begin{proof}
Since $p^*(\mathscr{F})$ and $q^*(\mathscr{G})$ are of finite type over
$\mathscr{O}_{X\times_S Y}$, we are reduced, by
(\hyperref[I.9.1.12]{9.1.12}) and
(\hyperref[0.1.7.5]{0, 1.7.5}), to the particular case where
$\mathscr{G}=\mathscr{O}_Y$, that is, to proving the formula
\[
  \operatorname{Supp}(p^{-1}(\mathscr{F}))
  =p^{-1}(\operatorname{Supp}(\mathscr{F})).
  \tag{9.1.13.1}\label{I.9.1.13.1}
\]

The same reasoning as in (\hyperref[0.1.7.5]{0, 1.7.5}) reduces us to
verifying that, for every $z\in X\times_S Y$, we have
$\mathscr{O}_z/\mathfrak{m}_x\mathscr{O}_z\neq0$ (with $x=p(z)$), which
follows from the fact that the homomorphism
$\mathscr{O}_x\to\mathscr{O}_z$ is \emph{local} by hypothesis.
\end{proof}

We leave it to the reader to extend the results proved in this subsection
for two factors to a product of an arbitrary number of factors.

\subsection{Direct image of a quasi-coherent sheaf.}
\label{subsection:I.9.2}

\begin{proposition}[9.2.1]
\label{I.9.2.1}
Let $f:X\to Y$ be a morphism of preschemes. Suppose that there exists a
covering $(Y_\alpha)$ of $Y$ by affine open subsets having the following
property: each of the $f^{-1}(Y_\alpha)$ admits a finite covering
$(X_{\alpha i})$ by affine open subsets contained in
$f^{-1}(Y_\alpha)$, such that each of the intersections
$X_{\alpha i}\cap X_{\alpha j}$ is itself a finite union of affine open
subsets. Under these conditions, for every quasi-coherent
$\mathscr{O}_X$-module $\mathscr{F}$, $f_*(\mathscr{F})$ is a
quasi-coherent $\mathscr{O}_Y$-module.
\end{proposition}

\begin{proof}
The question being local on $Y$, we may suppose $Y$ equal to one of the
$Y_\alpha$, and hence suppress the indices $\alpha$.

\begin{enumerate}
\item[a)] Suppose first that the $X_i\cap X_j$ are themselves
\emph{affine} open subsets. Set
$\mathscr{F}_i=\mathscr{F}|X_i$,
$\mathscr{F}_{ij}=\mathscr{F}|(X_i\cap X_j)$, and let
$\mathscr{F}'_i$ and $\mathscr{F}'_{ij}$ be the images of $\mathscr{F}_i$
and $\mathscr{F}_{ij}$ respectively under the restrictions of $f$ to $X_i$
and $X_i\cap X_j$; we know that the $\mathscr{F}'_i$ and the
$\mathscr{F}'_{ij}$ are quasi-coherent
(\hyperref[I.1.6.3]{1.6.3}). Set
$\mathscr{G}=\bigoplus_i\mathscr{F}'_i$,
$\mathscr{H}=\bigoplus_{i,j}\mathscr{F}'_{ij}$; $\mathscr{G}$ and
$\mathscr{H}$ are quasi-coherent $\mathscr{O}_Y$-modules; we will
define a homomorphism $u:\mathscr{G}\to\mathscr{H}$ such that
$f_*(\mathscr{F})$ is the \emph{kernel} of $u$; it will follow that
$f_*(\mathscr{F})$ is quasi-coherent
(\hyperref[I.1.3.9]{1.3.9}). It suffices to define $u$
\oldpage[I]{172}
as a homomorphism of presheaves; taking into account the definitions of
$\mathscr{G}$ and $\mathscr{H}$, it therefore suffices, for every open set
$W\subset Y$, to define a homomorphism
\[
  u_W:\bigoplus_i\Gamma(f^{-1}(W)\cap X_i,\mathscr{F})
  \longrightarrow
  \bigoplus_{i,j}\Gamma(f^{-1}(W)\cap X_i\cap X_j,\mathscr{F})
\]
so as to satisfy the usual compatibility conditions when $W$
varies. If, for every section
$s_i\in\Gamma(f^{-1}(W)\cap X_i,\mathscr{F})$, we denote by $s_{i|j}$ its
restriction to $f^{-1}(W)\cap X_i\cap X_j$, we set
\[
  u_W((s_i))=(s_{i|j}-s_{j|i})
\]
and the compatibility conditions are evidently satisfied. To
prove that the kernel $\mathscr{R}$ of $u$ is $f_*(\mathscr{F})$, let us
define a homomorphism of $f_*(\mathscr{F})$ into $\mathscr{R}$ by assigning
to every section $s\in\Gamma(f^{-1}(W),\mathscr{F})$ the family
$(s_i)$, where $s_i$ is the restriction of $s$ to $f^{-1}(W)\cap X_i$; the
axioms (F 1) and (F 2) of sheaves (G, II, 1.1) imply that this
homomorphism is \emph{bijective}, which proves the assertion in this case.

\item[b)] In the general case the same argument works as soon as we know that the $\mathscr{F}'_{ij}$ are quasi-coherent. Now, by
hypothesis, $X_i\cap X_j$ is a finite union of affine open subsets
$X_{ijk}$; and since the $X_{ijk}$ are affine open subsets \emph{in a
scheme}, the intersection of any two of them is again an affine open subset
(\hyperref[I.5.5.6]{5.5.6}). We are therefore reduced to the first case,
and (9.2.1) is thus proved.
\end{enumerate}
\end{proof}

\begin{corollary}[9.2.2]
\label{I.9.2.2}
The conclusion of (9.2.1) holds in each of the following cases:
\begin{enumerate}
\item[a)] $f$ is separated and quasi-compact.
\item[b)] $f$ is separated and of finite type.
\item[c)] $f$ is quasi-compact and the underlying space of $X$ is locally
Noetherian.
\end{enumerate}
\end{corollary}

\begin{proof}
In case a), the $X_{\alpha i}\cap X_{\alpha j}$ are affine
(\hyperref[I.5.5.6]{5.5.6}). Case b) is a particular case of a)
(\hyperref[I.6.6.3]{6.6.3}). Finally, in case c), we can reduce to the
case where $Y$ is affine and the underlying space of $X$ is Noetherian; then
$X$ admits a finite affine open covering $(X_i)$, and the $X_i\cap X_j$,
being quasi-compact, are finite unions of affine open subsets
(\hyperref[I.2.1.3]{2.1.3}).
\end{proof}

\subsection{Extension of sections of quasi-coherent sheaves.}
\label{subsection:I.9.3}

\begin{theorem}[9.3.1]
\label{I.9.3.1}
Let $X$ be a prescheme whose underlying space is Noetherian, or a scheme
whose underlying space is quasi-compact. Let $\mathscr{L}$ be an invertible
$\mathscr{O}_X$-module
(\hyperref[0.5.4.1]{0, 5.4.1}), $f$ a section of $\mathscr{L}$ over
$X$, $X_f$ the open set of the $x\in X$ where $f(x)\neq0$
(\hyperref[0.5.5.1]{0, 5.5.1}), and $\mathscr{F}$ a quasi-coherent
$\mathscr{O}_X$-module.
\begin{enumerate}
\item[(i)] If $s\in\Gamma(X,\mathscr{F})$ is such that $s|X_f=0$, there exists
an integer $n>0$ such that $s\otimes f^{\otimes n}=0$.
\item[(ii)] For every section $s\in\Gamma(X_f,\mathscr{F})$, there exists an
integer $n>0$ such that $s\otimes f^{\otimes n}$ extends to a section of
$\mathscr{F}\otimes\mathscr{L}^{\otimes n}$ over $X$.
\end{enumerate}
\end{theorem}

\begin{proof}
\begin{enumerate}
\item[(i)] Since the underlying space of $X$ is quasi-compact, hence a
finite union of affine open subsets $U_i$ such that $\mathscr{L}|U_i$ is
isomorphic to $\mathscr{O}_X|U_i$, we are reduced to the case where $X$ is
affine and $\mathscr{L}=\mathscr{O}_X$. In this case, $f$ is identified
with an element of $A(X)$ and we have $X_f=D(f)$; $s$ is identified with an
element of an $A(X)$-module $M$ and $s|X_f$ with the corresponding element
of $M_f$, and the result is trivial, taking into account the definition of
a module of fractions.
\oldpage[I]{173}

\item[(ii)] Again $X$ is a finite union of affine open subsets $U_i$
($1\leq i\leq r$) such that $\mathscr{L}|U_i\simeq\mathscr{O}_X|U_i$, and for
each $i$,
$(s\otimes f^{\otimes n})|(U_i\cap X_f)$ is identified, by the preceding
isomorphism, with
$(f|(U_i\cap X_f))^n(s|(U_i\cap X_f))$. We then know
(\hyperref[I.1.4.1]{1.4.1}) that there exists an integer $n>0$ such that
for each $i$, $(s\otimes f^{\otimes n})|(U_i\cap X_f)$ extends to a
section $s_i$ of $\mathscr{F}\otimes\mathscr{L}^{\otimes n}$ over
$U_i$. Let $s_{i|j}$ be the restriction of $s_i$ to $U_i\cap U_j$; by
definition we have $s_{i|j}-s_{j|i}=0$ in $X_f\cap U_i\cap U_j$. Now, if
$X$ is a Noetherian space, $U_i\cap U_j$ is quasi-compact; if $X$ is a
scheme, $U_i\cap U_j$ is an affine open subset
(\hyperref[I.5.5.6]{5.5.6}), hence again quasi-compact. By virtue of
(i), there therefore exists an integer $m$ (independent of $i$ and $j$)
such that $(s_{i|j}-s_{j|i})\otimes f^{\otimes m}=0$. We conclude
immediately that there exists a section $s'$ of
$\mathscr{F}\otimes\mathscr{L}^{\otimes(n+m)}$ over $X$, inducing
$s_i\otimes f^{\otimes m}$ over each $U_i$, and consequently inducing
$s\otimes f^{\otimes(n+m)}$ over $X_f$.
\end{enumerate}
\end{proof}

The corollaries that follow give an interpretation of Theorem (9.3.1) in a
more algebraic language:

\begin{corollary}[9.3.2]
\label{I.9.3.2}
The hypotheses being those of (9.3.1), consider the graded ring
$A_* = \Gamma_*(\mathscr{L})$ and the graded $A_*$-module
$M_* = \Gamma_*(\mathscr{L},\mathscr{F})$
(\hyperref[0.5.4.6]{0, 5.4.6}). If $f\in A_n$, where $n\in\mathbb{Z}$,
then we have a canonical isomorphism
\[
  \Gamma(X_f,\mathscr{F})\xrightarrow{\sim}((M_*)_f)_0
\]
(the subgroup of the module of fractions $(M_*)_f$ consisting of the
elements of degree $0$).
\end{corollary}

\begin{corollary}[9.3.3]
\label{I.9.3.3}
Suppose that the hypotheses of (9.3.1) are satisfied, and suppose moreover
that $\mathscr{L}=\mathscr{O}_X$. Then if we set
$A=\Gamma(X,\mathscr{O}_X)$, $M=\Gamma(X,\mathscr{F})$, the $A_f$-module
$\Gamma(X_f,\mathscr{F})$ is canonically isomorphic to $M_f$.
\end{corollary}

\begin{proposition}[9.3.4]
\label{I.9.3.4}
Let $X$ be a Noetherian prescheme, $\mathscr{F}$ a coherent
$\mathscr{O}_X$-module, and $\mathscr{J}$ a coherent sheaf of ideals in
$\mathscr{O}_X$ such that the support of $\mathscr{F}$ is contained in
that of $\mathscr{O}_X/\mathscr{J}$. Then there exists an integer $n>0$
such that $\mathscr{J}^n\mathscr{F}=0$.
\end{proposition}

\begin{proof}
Since $X$ is a finite union of affine open subsets whose rings are
Noetherian, we may suppose $X$ affine with Noetherian ring $A$; then
$\mathscr{F}=\widetilde M$, where $M=\Gamma(X,\mathscr{F})$ is an $A$-module
of finite type, and $\mathscr{J}=\widetilde{\mathfrak{J}}$, where
$\mathfrak{J}=\Gamma(X,\mathscr{J})$ is an ideal of $A$
(\hyperref[I.1.4.1]{1.4.1} and \hyperref[I.1.5.1]{1.5.1}). Since $A$
is Noetherian, $\mathfrak{J}$ admits a finite system of generators $f_i$
($1\leq i\leq m$). By hypothesis, every section of $\mathscr{F}$ over $X$
is zero in each of the $D(f_i)$; if $s_j$ ($1\leq j\leq q$) are sections of
$\mathscr{F}$ generating $M$, there therefore exists an integer $h$,
independent of $i$ and $j$, such that $f_i^h s_j=0$
(\hyperref[I.1.4.1]{1.4.1}), hence $f_i^h s=0$ for every $s\in M$. We
conclude that if $n=mh$, we have $\mathfrak{J}^nM=0$, and consequently the
corresponding $\mathscr{O}_X$-module
$\mathscr{J}^n\mathscr{F}=\widetilde{\mathfrak{J}^nM}$
(\hyperref[I.1.3.13]{1.3.13}) is zero.
\end{proof}

\begin{corollary}[9.3.5]
\label{I.9.3.5}
Under the hypotheses of (9.3.4), there exists a closed sub-prescheme $Y$ of
$X$ whose underlying space is the support of
$\mathscr{O}_X/\mathscr{J}$, and such that, if $j:Y\to X$ is the canonical
injection, we have $\mathscr{F}=j_*(j^*(\mathscr{F}))$.
\end{corollary}

\begin{proof}
Let us note first that the supports of $\mathscr{O}_X/\mathscr{J}$ and of
$\mathscr{O}_X/\mathscr{J}^n$ are the same, since if
$\mathscr{J}_x=\mathscr{O}_x$, we also have
$\mathscr{J}_x^n=\mathscr{O}_x$, and on the other hand we have
$\mathscr{J}_x^n\subset\mathscr{J}_x$ for every $x\in X$. We may therefore,
by virtue of (9.3.4), suppose that $\mathscr{J}\mathscr{F}=0$; we then take
for $Y$ the closed sub-prescheme of $X$ defined by $\mathscr{J}$, and since
$\mathscr{F}$ is then an $(\mathscr{O}_X/\mathscr{J})$-module, the
conclusion is immediate.
\end{proof}

\oldpage[I]{174}

\subsection{Extension of quasi-coherent sheaves.}
\label{subsection:I.9.4}

\begin{env}[9.4.1]
\label{I.9.4.1}
Let $X$ be a topological space, $\mathscr{F}$ a sheaf of sets
(resp. of groups, of rings) on $X$, $U$ an open subset of $X$,
$\psi:U\to X$ the canonical injection, and $\mathscr{G}$ a sub-sheaf of
$\mathscr{F}|U=\psi^*(\mathscr{F})$. Since $\psi_*$ is left exact,
$\psi_*(\mathscr{G})$ is a sub-sheaf of
$\psi_*(\psi^*(\mathscr{F}))$; if we consider the canonical homomorphism
$\rho:\mathscr{F}\to\psi_*(\psi^*(\mathscr{F}))$
(\hyperref[0.3.5.3]{0, 3.5.3}), we will denote by
$\overline{\mathscr{G}}$ the sub-sheaf
$\rho^{-1}(\psi_*(\mathscr{G}))$ of $\mathscr{F}$. It follows immediately
from the definitions that for every open subset $V$ of $X$,
$\Gamma(V,\overline{\mathscr{G}})$ consists of the sections
$s\in\Gamma(V,\mathscr{F})$ whose restriction to $V\cap U$ is a section
of $\mathscr{G}$ over $V\cap U$. We therefore have
$\overline{\mathscr{G}}|U=\psi^*(\overline{\mathscr{G}})=\mathscr{G}$, and
$\overline{\mathscr{G}}$ is the \emph{largest} sub-sheaf of
$\mathscr{F}$ inducing $\mathscr{G}$ on $U$; we will say that
$\overline{\mathscr{G}}$ is the \emph{canonical extension} of the
sub-sheaf $\mathscr{G}$ of $\mathscr{F}|U$ to a sub-sheaf of
$\mathscr{F}$.
\end{env}

\begin{proposition}[9.4.2]
\label{I.9.4.2}
Let $X$ be a prescheme and $U$ an open subset of $X$ such that the canonical
injection $j:U\to X$ is a quasi-compact morphism
\emph{(which will be satisfied for \emph{every} $U$ if the underlying space
$X$ is \emph{locally Noetherian}
(\hyperref[I.6.6.4]{6.6.4, (i)}))}. Then:
\begin{enumerate}
\item[(i)] For every quasi-coherent $(\mathscr{O}_X|U)$-module
  $\mathscr{G}$, $j_*(\mathscr{G})$ is a quasi-coherent
  $\mathscr{O}_X$-module and we have
  $j_*(\mathscr{G})|U=j^*(j_*(\mathscr{G}))=\mathscr{G}$.
\item[(ii)] For every quasi-coherent $\mathscr{O}_X$-module $\mathscr{F}$
  and every quasi-coherent sub-$(\mathscr{O}_X|U)$-module $\mathscr{G}$, the
  canonical extension $\overline{\mathscr{G}}$ of $\mathscr{G}$ (9.4.1)
  is a quasi-coherent sub-$\mathscr{O}_X$-module of $\mathscr{F}$.
\end{enumerate}
\end{proposition}

\begin{proof}
If $j=(\psi,\theta)$ ($\psi$ being the injection $U\to X$ of the
underlying spaces), we have by definition $j_*(\mathscr{G})=\psi_*(\mathscr{G})$
for every $(\mathscr{O}_X|U)$-module $\mathscr{G}$, and moreover
$j^*(\mathscr{H})=\psi^*(\mathscr{H})=\mathscr{H}|U$ for every
$\mathscr{O}_X$-module $\mathscr{H}$, by reason of the definition of a
prescheme induced on an open subset. (i) is therefore a particular case of
(\hyperref[I.9.2.2]{9.2.2, a)}); for the same reason,
$j_*(j^*(\mathscr{F}))$ is quasi-coherent, and since
$\overline{\mathscr{G}}$ is the inverse image of $j_*(\mathscr{G})$ under
the homomorphism $\rho:\mathscr{F}\to j_*(j^*(\mathscr{F}))$, (ii) follows
from (\hyperref[I.4.1.1]{4.1.1}).
\end{proof}

Note that the hypothesis that the morphism $j:U\to X$ is quasi-compact is
also satisfied when the open subset $U$ is \emph{quasi-compact} and $X$ is a
\emph{scheme}: indeed, $U$ is then a finite union of affine open subsets
$U_i$, and for every affine open subset $V$ of $X$, $V\cap U_i$ is an affine
open subset (\hyperref[I.5.5.6]{5.5.6}), hence quasi-compact.

\begin{corollary}[9.4.3]
\label{I.9.4.3}
Let $X$ be a prescheme and $U$ a quasi-compact open subset of $X$ such that
the injection morphism $j:U\to X$ is quasi-compact. Suppose moreover that
every quasi-coherent $\mathscr{O}_X$-module is the inductive limit of its
quasi-coherent sub-$\mathscr{O}_X$-modules of finite type
\emph{(which is the case when $X$ is an \emph{affine scheme})}. Let then
$\mathscr{F}$ be a quasi-coherent $\mathscr{O}_X$-module, and $\mathscr{G}$
a quasi-coherent sub-$(\mathscr{O}_X|U)$-module of finite type of
$\mathscr{F}|U$. Then there exists a quasi-coherent
sub-$\mathscr{O}_X$-module of finite type $\mathscr{G}'$ of $\mathscr{F}$
such that $\mathscr{G}'|U=\mathscr{G}$.
\end{corollary}

\begin{proof}
Indeed, we have $\mathscr{G}=\overline{\mathscr{G}}|U$, and
$\overline{\mathscr{G}}$ is quasi-coherent by (9.4.2), hence the inductive
limit of its quasi-coherent sub-$\mathscr{O}_X$-modules of finite type
$\mathscr{H}_\lambda$. Consequently $\mathscr{G}$ is the inductive limit of
the $\mathscr{H}_\lambda|U$, hence equal to one of the
$\mathscr{H}_\lambda|U$ since it is of finite type
(\hyperref[0.5.2.3]{0, 5.2.3}).
\end{proof}

\begin{remark}[9.4.4]
\label{I.9.4.4}
Suppose that for \emph{every} affine open subset $U\subset X$ the injection
morphism $U\to X$ is quasi-compact. Then if the conclusion of (9.4.3)
holds for every affine open subset $U$ and every quasi-coherent
sub-$(\mathscr{O}_X|U)$-module of finite type $\mathscr{G}$ of
$\mathscr{F}|U$,
\oldpage[I]{175}
it follows that $\mathscr{F}$ is the inductive limit of its quasi-coherent
sub-$\mathscr{O}_X$-modules of finite type. Indeed, for every affine open
subset $U\subset X$, we have $\mathscr{F}|U=\widetilde{M}$, where $M$ is an
$A(U)$-module, and since the latter is the inductive limit of its
submodules of finite type, $\mathscr{F}|U$ is the inductive limit of its
quasi-coherent sub-$(\mathscr{O}_X|U)$-modules of finite type
(\hyperref[I.1.3.9]{1.3.9}). Now, by hypothesis, each of these
submodules is induced on $U$ by a quasi-coherent sub-$\mathscr{O}_X$-module
of finite type $\mathscr{G}_{\lambda,U}$ of $\mathscr{F}$. The finite sums
of the $\mathscr{G}_{\lambda,U}$ are again quasi-coherent
$\mathscr{O}_X$-modules of finite type, since the question is local and the
case where $X$ is affine has been treated in
(\hyperref[I.1.3.10]{1.3.10}); it is then clear that $\mathscr{F}$ is the
inductive limit of these finite sums, whence our assertion.
\end{remark}

\begin{corollary}[9.4.5]
\label{I.9.4.5}
Under the hypotheses of (9.4.3), for every quasi-coherent
$(\mathscr{O}_X|U)$-module of finite type $\mathscr{G}$, there exists a
quasi-coherent $\mathscr{O}_X$-module of finite type $\mathscr{G}'$ such
that $\mathscr{G}'|U=\mathscr{G}$.
\end{corollary}

\begin{proof}
Since $\mathscr{F}=j_*(\mathscr{G})$ is quasi-coherent (9.4.2) and
$\mathscr{F}|U=\mathscr{G}$, it suffices to apply (9.4.3) to $\mathscr{F}$.
\end{proof}

\begin{lemma}[9.4.6]
\label{I.9.4.6}
Let $X$ be a prescheme, $L$ a well-ordered set,
$(V_\lambda)_{\lambda\in L}$ a covering of $X$ by affine open subsets, and
$U$ an open subset of $X$; for every $\lambda\in L$, set
$W_\lambda=\bigcup_{\mu<\lambda}V_\mu$. Suppose that:
\begin{enumerate}
  \item[1\textsuperscript{o}] For every $\lambda\in L$,
    $V_\lambda\cap W_\lambda$ is quasi-compact;
  \item[2\textsuperscript{o}] The immersion morphism $U\to X$ is
    quasi-compact.
\end{enumerate}
Then, for every quasi-coherent $\mathscr{O}_X$-module $\mathscr{F}$ and every
quasi-coherent sub-$(\mathscr{O}_X|U)$-module of finite type $\mathscr{G}$ of
$\mathscr{F}|U$, there exists a quasi-coherent sub-$\mathscr{O}_X$-module of
finite type $\mathscr{G}'$ of $\mathscr{F}$ such that
$\mathscr{G}'|U=\mathscr{G}$.
\end{lemma}

\begin{proof}
Let $\kappa$ be the order type of $L$; let us index the given covering by $\lambda<\kappa$.\footnote{Editorial clarification: the printed induction omits the initial case and, if $L$ has a largest element, the final step. The steps $0\leq\lambda\leq\kappa$ complete the proof without changing the hypotheses.}
For $0\leq\lambda\leq\kappa$, set $W_\lambda=\bigcup_{\mu<\lambda}V_\mu$ and $U_\lambda=U\cup W_\lambda$.
We construct by transfinite induction a quasi-coherent sub-$(\mathscr{O}_X|U_\lambda)$-module of finite type $\mathscr{G}'_\lambda$ of $\mathscr{F}|U_\lambda$ such that $\mathscr{G}'_\lambda|U_\mu=\mathscr{G}'_\mu$ for $\mu<\lambda$ and $\mathscr{G}'_\lambda|U=\mathscr{G}$.
We begin with $\mathscr{G}'_0=\mathscr{G}$ on $U_0=U$; since $U_\kappa=X$, the module $\mathscr{G}'=\mathscr{G}'_\kappa$ will answer the question.
Suppose therefore that $0<\lambda\leq\kappa$ and that the $\mathscr{G}'_\mu$ have been constructed with these properties for $\mu<\lambda$.
If $\lambda$ is a limit ordinal, we take for $\mathscr{G}'_\lambda$ the unique sub-$(\mathscr{O}_X|U_\lambda)$-module of $\mathscr{F}|U_\lambda$ such that $\mathscr{G}'_\lambda|U_\mu=\mathscr{G}'_\mu$ for every $\mu<\lambda$; the $U_\mu$ with $\mu<\lambda$ then cover $U_\lambda$, and quasi-coherence and finite type are local properties.
If, on the contrary, $\lambda=\mu+1$, we have $U_\lambda=U_\mu\cup V_\mu$, and it will suffice to define a
quasi-coherent sub-$(\mathscr{O}_X|V_\mu)$-module of finite type
$\mathscr{G}''_\mu$ of $\mathscr{F}|V_\mu$ such that
\[
  \mathscr{G}''_\mu|(U_\mu\cap V_\mu)
  =\mathscr{G}'_\mu|(U_\mu\cap V_\mu)~;
\]
we will then take for $\mathscr{G}'_\lambda$ the
sub-$(\mathscr{O}_X|U_\lambda)$-module of $\mathscr{F}|U_\lambda$ such that
$\mathscr{G}'_\lambda|U_\mu=\mathscr{G}'_\mu$ and
$\mathscr{G}'_\lambda|V_\mu=\mathscr{G}''_\mu$
(\hyperref[0.3.3.1]{0, 3.3.1}). Now, since $V_\mu$ is affine, the
existence of $\mathscr{G}''_\mu$ is ensured by (9.4.3) as soon as we have
proved that $U_\mu\cap V_\mu$ is quasi-compact; but $U_\mu\cap V_\mu$ is the
union of $U\cap V_\mu$ and $W_\mu\cap V_\mu$, both of which are
quasi-compact by virtue of the hypothesis.
\end{proof}

\begin{theorem}[9.4.7]
\label{I.9.4.7}
Let $X$ be a prescheme and $U$ an open subset of $X$. Suppose that one of
the following conditions is satisfied:
\begin{enumerate}
  \item[(a)] The underlying space of $X$ is locally Noetherian.
  \item[(b)] $X$ is a quasi-compact scheme and $U$ a quasi-compact open
  subset.
\end{enumerate}
Then for every quasi-coherent $\mathscr{O}_X$-module $\mathscr{F}$ and every
quasi-coherent sub-$(\mathscr{O}_X|U)$-module of finite type $\mathscr{G}$ of
$\mathscr{F}|U$, there exists a quasi-coherent sub-$\mathscr{O}_X$-module
of finite type $\mathscr{G}'$ of $\mathscr{F}$ such that
$\mathscr{G}'|U=\mathscr{G}$.
\end{theorem}

\oldpage[I]{176}

\begin{proof}
Let $(V_\lambda)_{\lambda\in L}$ be a covering of $X$ by affine open
subsets, $L$ being supposed finite in case b); $L$ being endowed with a
structure of well-ordered set, it suffices to verify that the conditions
of Lemma (9.4.6) are satisfied. This is evident under hypothesis a), the
spaces $V_\lambda$ being Noetherian. Under hypothesis b), the
$V_\lambda\cap V_\mu$ are affine (\hyperref[I.5.5.6]{5.5.6}), hence
quasi-compact, and since $L$ is finite, $V_\lambda\cap W_\lambda$ is
quasi-compact. Whence the theorem.
\end{proof}

\begin{corollary}[9.4.8]
\label{I.9.4.8}
Under the hypotheses of (9.4.7), for every quasi-coherent
$(\mathscr{O}_X|U)$-module of finite type $\mathscr{G}$, there exists a
quasi-coherent $\mathscr{O}_X$-module of finite type $\mathscr{G}'$ such
that $\mathscr{G}'|U=\mathscr{G}$.
\end{corollary}

\begin{proof}
It suffices to apply (9.4.7) to $\mathscr{F}=j_*(\mathscr{G})$, which is
quasi-coherent (9.4.2) and such that $\mathscr{F}|U=\mathscr{G}$.
\end{proof}

\begin{corollary}[9.4.9]
\label{I.9.4.9}
Let $X$ be a prescheme whose underlying space is locally Noetherian, or
a quasi-compact scheme. Then every quasi-coherent $\mathscr{O}_X$-module is
the inductive limit of its quasi-coherent sub-$\mathscr{O}_X$-modules of
finite type.
\end{corollary}

\begin{proof}
This follows from (9.4.7) and from Remark (9.4.4).
\end{proof}

\begin{corollary}[9.4.10]
\label{I.9.4.10}
Under the hypotheses of (9.4.9), if a quasi-coherent $\mathscr{O}_X$-module
$\mathscr{F}$ is such that every quasi-coherent sub-$\mathscr{O}_X$-module
of finite type of $\mathscr{F}$ is generated by its sections over $X$,
then $\mathscr{F}$ is generated by its sections over $X$.
\end{corollary}

\begin{proof}
Indeed, let $U$ be an affine open neighbourhood of a point $x\in X$, and let
$s$ be a section of $\mathscr{F}$ over $U$; the
sub-$\mathscr{O}_X$-module $\mathscr{G}$ of $\mathscr{F}|U$ generated by
$s$ is quasi-coherent and of finite type, so there exists a
quasi-coherent sub-$\mathscr{O}_X$-module of finite type $\mathscr{G}'$ of
$\mathscr{F}$ such that $\mathscr{G}'|U=\mathscr{G}$ (9.4.7). By hypothesis,
there are therefore finitely many sections $t_i$ of $\mathscr{G}'$ over
$X$ and sections $a_i$ of $\mathscr{O}_X$ over a neighbourhood
$V\subset U$ of $x$ such that
$s|V=\sum_i a_i\mathbin{\cdot}(t_i|V)$, which proves the corollary.
\end{proof}


\subsection{Closed image of a prescheme. Closure of a sub-prescheme.}
\label{subsection:I.9.5}

\begin{proposition}[9.5.1]
\label{I.9.5.1}
Let $f:X\to Y$ be a morphism of preschemes such that $f_*(\mathscr{O}_X)$ is
a quasi-coherent $\mathscr{O}_Y$-module (which is the case if $f$ is
quasi-compact and if moreover $f$ is separated or $X$ is locally Noetherian
(\hyperref[I.9.2.2]{9.2.2})). Then there exists a smallest closed
sub-prescheme $Y'$ of $Y$ such that the canonical injection $j:Y'\to Y$
dominates $f$ (or, what amounts to the same (\hyperref[I.4.4.1]{4.4.1}),
such that the sub-prescheme $f^{-1}(Y')$ of $X$ is \emph{identical} to $X$).
\end{proposition}

More precisely:

\begin{corollary}[9.5.2]
\label{I.9.5.2}
Under the conditions of (9.5.1), let $f=(\psi,\theta)$ and let $\mathscr{J}$
be the (quasi-coherent) kernel of the homomorphism
$\theta:\mathscr{O}_Y\to f_*(\mathscr{O}_X)$. Then the closed sub-prescheme
$Y'$ of $Y$ defined by $\mathscr{J}$ satisfies the conditions of (9.5.1).
\end{corollary}

\begin{proof}
Since the functor $\psi^*$ is exact, the canonical factorization
$\theta:\mathscr{O}_Y\to\mathscr{O}_Y/\mathscr{J}
\xrightarrow{\theta'}\psi_*(\mathscr{O}_X)$ gives
(\hyperref[0.3.5.4.3]{0, 3.5.4.3}) a factorization
\[
  \theta^\sharp:\psi^*(\mathscr{O}_Y)
  \to\psi^*(\mathscr{O}_Y)/\psi^*(\mathscr{J})
  \xrightarrow{{\theta'}^\sharp}\mathscr{O}_X~;
\]
since for every $x\in X$, $\theta_x^\sharp$ is a local homomorphism, the
same is true of ${\theta'_x}^\sharp$; if we denote by $\psi_0$ the
continuous map $\psi$ considered as a map from $X$ to $X'$, and by
$\theta_0$ the restriction
$\theta'|X':(\mathscr{O}_Y/\mathscr{J})|X'
\to\psi_*(\mathscr{O}_X)|X'=(\psi_0)_*(\mathscr{O}_X)$, we therefore see
that $f_0=(\psi_0,\theta_0)$ is a morphism of preschemes $X\to X'$
(\hyperref[I.2.2.1]{2.2.1}) such that $f=j\circ f_0$. If now $X''$
\oldpage[I]{177}
is a second closed sub-prescheme of $Y$, defined by a quasi-coherent sheaf
of ideals $\mathscr{J}'$ of $\mathscr{O}_Y$, and such that the injection
$j':X''\to Y$ dominates $f$, we must first of all have
$X''\supset\psi(X)$, hence $X'\subset X''$ since $X''$ is closed. Moreover,
for every $y\in X''$, $\theta$ must factor as
$\mathscr{O}_y\to\mathscr{O}_y/\mathscr{J}'_y
\to(\psi_*(\mathscr{O}_X))_y$, which by definition implies
$\mathscr{J}'_y\subset\mathscr{J}_y$, and consequently $X'$ is a closed
sub-prescheme of $X''$ (\hyperref[I.4.1.10]{4.1.10}).
\end{proof}

\begin{definition}[9.5.3]
\label{I.9.5.3}
When there exists a smallest closed sub-prescheme $Y'$ of $Y$ such that
the canonical injection $j:Y'\to Y$ dominates $f$, we say that $Y'$ is the
closed image prescheme of $X$ under the morphism $f$.
\end{definition}

\begin{proposition}[9.5.4]
\label{I.9.5.4}
If $f_*(\mathscr{O}_X)$ is a quasi-coherent $\mathscr{O}_Y$-module, the
underlying space of the closed image of $X$ under $f$ is the closure
$\overline{f(X)}$ in $Y$.
\end{proposition}

\begin{proof}
Since the support of $f_*(\mathscr{O}_X)$ is contained in
$\overline{f(X)}$, we have (with the notation of
(\hyperref[I.9.5.2]{9.5.2})) $\mathscr{J}_y=\mathscr{O}_y$ for
$y\notin\overline{f(X)}$, hence the support of
$\mathscr{O}_Y/\mathscr{J}$ is contained in $\overline{f(X)}$. Moreover,
this support is closed and contains $f(X)$: indeed, if $y\in f(X)$, the
unit element of the ring $(\psi_*(\mathscr{O}_X))_y$ is not zero, being the
germ at $y$ of the section
\[
  1\in\Gamma(X,\mathscr{O}_X)=\Gamma(Y,\psi_*(\mathscr{O}_X))~;
\]
since it is the image under $\theta$ of the unit element of $\mathscr{O}_y$,
the latter does not belong to $\mathscr{J}_y$, hence
$\mathscr{O}_y/\mathscr{J}_y\neq0$; this completes the proof.
\end{proof}

\begin{proposition}[9.5.5]
\label{I.9.5.5}
\emph{(transitivity of closed images).} Let $f:X\to Y$ and $g:Y\to Z$
be two morphisms of preschemes; suppose that the closed image $Y'$ of $X$
under $f$ exists, and that, if $g'$ is the restriction of $g$ to $Y'$, the
closed image $Z'$ of $Y'$ under $g'$ exists. Then the closed image of $X$
under $g\circ f$ exists and is equal to $Z'$.
\end{proposition}

\begin{proof}
It suffices (\hyperref[I.9.5.1]{9.5.1}) to show that $Z'$ is the
smallest closed sub-prescheme $Z_1$ of $Z$ such that the closed sub-prescheme
$(g\circ f)^{-1}(Z_1)$ of $X$ (equal to $f^{-1}(g^{-1}(Z_1))$ by
(\hyperref[I.4.4.2]{4.4.2})) is equal to $X$; it amounts to the same to
say that $Z'$ is the smallest closed sub-prescheme of $Z$ such that $f$ is
dominated by the injection $g^{-1}(Z_1)\to Y$
(\hyperref[I.4.4.1]{4.4.1}). Now, by virtue of the existence of the
closed image $Y'$, every $Z_1$ having this property is such that
$g^{-1}(Z_1)$ dominates $Y'$, which is equivalent to saying that
$j^{-1}(g^{-1}(Z_1))=g'^{-1}(Z_1)=Y'$, denoting by $j$ the injection
$Y'\to Y$. By definition of $Z'$, we indeed conclude that $Z'$ is the
smallest closed sub-prescheme of $Z$ satisfying the preceding condition.
\end{proof}

\begin{corollary}[9.5.6]
\label{I.9.5.6}
Let $f:X\to Y$ be an $S$-morphism such that $Y$ is the closed image of $X$
under $f$. Let $Z$ be an $S$-scheme; if two $S$-morphisms $g_1$, $g_2$ from
$Y$ to $Z$ are such that $g_1\circ f=g_2\circ f$, then $g_1=g_2$.
\end{corollary}

\begin{proof}
Let $h=(g_1,g_2)_S:Y\to Z\times_S Z$; since the diagonal
$T=\Delta_Z(Z)$ is a closed sub-prescheme of $Z\times_S Z$,
$Y'=h^{-1}(T)$ is a closed sub-prescheme of $Y$
(\hyperref[I.4.4.1]{4.4.1}). Set $u=g_1\circ f=g_2\circ f$; we then
have, by definition of the product, $h'=h\circ f=(u,u)_S$, hence
$h\circ f=\Delta_Z\circ u$; since $\Delta_Z^{-1}(T)=Z$, we have
$h'^{-1}(T)=u^{-1}(Z)=X$, hence $f^{-1}(Y')=X$. We conclude
(\hyperref[I.4.4.1]{4.4.1}) that the canonical injection $Y'\to Y$
dominates $f$, hence $Y'=Y$ by hypothesis; consequently
(\hyperref[I.4.4.1]{4.4.1}), $h$ factors as $\Delta_Z\circ v$, where
$v$ is a morphism $Y\to Z$, which implies $g_1=g_2=v$.
\end{proof}

\begin{remark}[9.5.7]
\label{I.9.5.7}
If $X$ and $Y$ are $S$-schemes, Proposition
(\hyperref[I.9.5.6]{9.5.6}) means that
\oldpage[I]{178}
when $Y$ is the closed image of $X$ under $f$, $f$ is an
\emph{epimorphism} in the category of $S$-schemes (T, 1.1). We will prove
in Chapter V that, conversely, if the closed image $Y'$ of $X$ under $f$
exists and if $f$ is an epimorphism of $S$-schemes, then we necessarily
have $Y'=Y$.
\end{remark}

\begin{proposition}[9.5.8]
\label{I.9.5.8}
Suppose that the hypotheses of (\hyperref[I.9.5.1]{9.5.1}) are satisfied,
and let $Y'$ be the closed image of $X$ under $f$. For every open subset $V$
of $Y$, let $f_V:f^{-1}(V)\to V$ be the restriction of $f$; then the closed
image of $f^{-1}(V)$ under $f_V$ in $V$ exists and is equal to the
prescheme induced by $Y'$ on the open subset $V\cap Y'$ of $Y'$ (in other
words, to the sub-prescheme $\inf(V,Y')$ of $Y$
(\hyperref[I.4.4.3]{4.4.3})).
\end{proposition}

\begin{proof}
Set $X'=f^{-1}(V)$; since the direct image of $\mathscr{O}_{X'}$ under
$f_V$ is none other than the restriction of $f_*(\mathscr{O}_X)$ to $V$, it
is clear that the kernel $\mathscr{J}'$ of the homomorphism
$\mathscr{O}_V\to(f_V)_*(\mathscr{O}_{X'})$ is the restriction of
$\mathscr{J}$ to $V$, whence the proposition immediately.
\end{proof}

Note that this result can be interpreted by saying that the formation of
the closed image commutes with an extension $Y_1\to Y$ of the base
prescheme which is an \emph{open immersion}. We will see in Chapter IV that
the same holds for an extension $Y_1\to Y$ which is a \emph{flat} morphism,
provided that $f$ is separated and quasi-compact.

\begin{proposition}[9.5.9]
\label{I.9.5.9}
Let $f:X\to Y$ be a morphism such that the closed image $Y'$ of $X$ under
$f$ exists.
\begin{enumerate}
  \item[(i)] If $X$ is reduced, so is $Y'$.
  \item[(ii)] If we suppose that the hypotheses of
    (\hyperref[I.9.5.1]{9.5.1}) are satisfied and if $X$ is irreducible
    (resp. integral), so is $Y'$.
\end{enumerate}
\end{proposition}

\begin{proof}
By hypothesis, the morphism $f$ factors as
$X\xrightarrow{g}Y'\xrightarrow{j}Y$, where $j$ is the canonical injection.
Since $X$ is reduced, $g$ factors as
$X\xrightarrow{h}Y'_{\mathrm{red}}\xrightarrow{j'}Y'$, where $j'$ is
the canonical injection (\hyperref[I.5.2.2]{5.2.2}), and it then follows
from the definition of $Y'$ that $Y'_{\mathrm{red}}=Y'$. If moreover the
conditions of (\hyperref[I.9.5.1]{9.5.1}) are fulfilled, it follows from
(\hyperref[I.9.5.4]{9.5.4}) that $f(X)$ is dense in $Y'$; if $X$ is
irreducible, the same is therefore true of $Y'$
(\hyperref[0.2.1.5]{0, 2.1.5}). The assertion concerning integral
preschemes follows from the conjunction of the other two.
\end{proof}

\begin{proposition}[9.5.10]
\label{I.9.5.10}
Let $Y$ be a sub-prescheme of a prescheme $X$ such that the canonical
injection $i:Y\to X$ is a quasi-compact morphism. Then there exists a
smallest closed sub-prescheme $\overline{Y}$ of $X$ dominating $Y$; its
underlying space is the closure of that of $Y$; the latter is open in its
closure, and the prescheme $Y$ is induced on this open subset by
$\overline{Y}$.
\end{proposition}

\begin{proof}
It suffices to apply (\hyperref[I.9.5.1]{9.5.1}) to the injection $j$,
which is separated (\hyperref[I.5.5.1]{5.5.1}) and quasi-compact by
hypothesis; (\hyperref[I.9.5.1]{9.5.1}) therefore proves the existence of
$\overline{Y}$ and (\hyperref[I.9.5.4]{9.5.4}) shows that its underlying
space is the closure of $Y$ in $X$; since $Y$ is locally closed in $X$, it
is open in $\overline{Y}$, and the last assertion comes from
(\hyperref[I.9.5.8]{9.5.8}) applied to an open subset $V$ of $X$ such
that $Y$ is closed in $V$.
\end{proof}

With this notation, if the injection $V\to X$ is quasi-compact, and if
$\mathscr{J}$ is the quasi-coherent sheaf of ideals of
$\mathscr{O}_X|V$ defining the closed sub-prescheme $Y$ of $V$, it follows
from (\hyperref[I.9.5.1]{9.5.1}) that the quasi-coherent sheaf of ideals
of $\mathscr{O}_X$ defining $\overline{Y}$ is the canonical extension
(\hyperref[I.9.4.1]{9.4.1}) $\overline{\mathscr{J}}$ of $\mathscr{J}$,
since it is evidently the largest quasi-coherent sub-sheaf of ideals of
$\mathscr{O}_X$ inducing $\mathscr{J}$ on $V$.
\oldpage[I]{179}
\begin{corollary}[9.5.11]
\label{I.9.5.11}
Under the hypotheses of (\hyperref[I.9.5.10]{9.5.10}), every section of
$\mathscr{O}_{\overline{Y}}$ over an open subset $V$ of $\overline{Y}$
which is zero in $V\cap Y$ is zero.
\end{corollary}

\begin{proof}
By virtue of (\hyperref[I.9.5.8]{9.5.8}), we can reduce to the case where
$V=\overline{Y}$. If we take into account that the sections of
$\mathscr{O}_{\overline{Y}}$ over $\overline{Y}$ correspond
canonically to the $\overline{Y}$-sections of
$\overline{Y}\otimes_{\mathbb{Z}}\mathbb{Z}[T]$
(\hyperref[I.3.3.15]{3.3.15}) and that the latter is separated over
$\overline{Y}$, the corollary appears as a particular case of
(\hyperref[I.9.5.6]{9.5.6}).
\end{proof}

When there exists a smallest closed sub-prescheme $Y'$ of $X$
dominating a sub-prescheme $Y$ of $X$, we say that $Y'$ is the
\emph{closure} of $Y$ in $X$, when no confusion can result.

\subsection{Quasi-coherent sheaves of algebras; change of the structural
sheaf.}
\label{subsection:I.9.6}

\begin{proposition}[9.6.1]
\label{I.9.6.1}
Let $X$ be a prescheme and $\mathscr{B}$ a quasi-coherent
$\mathscr{O}_X$-algebra (\hyperref[0.5.1.3]{0, 5.1.3}). For a
$\mathscr{B}$-module $\mathscr{F}$ to be quasi-coherent (on the ringed
space $(X,\mathscr{B})$), it is necessary and sufficient that $\mathscr{F}$
be a quasi-coherent $\mathscr{O}_X$-module.
\end{proposition}

\begin{proof}
Since the question is local, we may suppose $X$ affine with ring $A$, and
then $\mathscr{B}=\widetilde{B}$, where $B$ is an $A$-algebra
(\hyperref[I.1.4.3]{1.4.3}). If $\mathscr{F}$ is quasi-coherent on
the ringed space $(X,\mathscr{B})$, we may also suppose that
$\mathscr{F}$ is the cokernel of a $\mathscr{B}$-homomorphism
$\mathscr{B}^{(I)}\to\mathscr{B}^{(J)}$; since this homomorphism is
also an $\mathscr{O}_X$-homomorphism of $\mathscr{O}_X$-modules, and
$\mathscr{B}^{(I)}$ and $\mathscr{B}^{(J)}$ are quasi-coherent
$\mathscr{O}_X$-modules
(\hyperref[I.1.3.9]{1.3.9, (ii)}), $\mathscr{F}$ is also a
quasi-coherent $\mathscr{O}_X$-module
(\hyperref[I.1.3.9]{1.3.9, (i)}).

Conversely, if $\mathscr{F}$ is a quasi-coherent $\mathscr{O}_X$-module,
we have $\mathscr{F}=\widetilde{M}$, where $M$ is a $B$-module
(\hyperref[I.1.4.3]{1.4.3}); $M$ is isomorphic to a cokernel of a
$B$-homomorphism $B^{(I)}\to B^{(J)}$, hence $\mathscr{F}$ is a
$\mathscr{B}$-module isomorphic to the cokernel of the corresponding
homomorphism $\mathscr{B}^{(I)}\to\mathscr{B}^{(J)}$
(\hyperref[I.1.3.13]{1.3.13}), which completes the proof.
\end{proof}

In particular, if $\mathscr{F}$ and $\mathscr{G}$ are two quasi-coherent
$\mathscr{B}$-modules,
$\mathscr{F}\otimes_{\mathscr{B}}\mathscr{G}$ is a quasi-coherent
$\mathscr{B}$-module; the same is true of
$\mathscr{H}\!om_{\mathscr{B}}(\mathscr{F},\mathscr{G})$ when we suppose
moreover that $\mathscr{F}$ admits a finite presentation
(\hyperref[I.1.3.13]{1.3.13}).

\begin{env}[9.6.2]
\label{I.9.6.2}
Given a prescheme $X$, we will say that a quasi-coherent
$\mathscr{O}_X$-algebra $\mathscr{B}$ is of \emph{finite type} if for every
$x\in X$ there exists an \emph{affine} open neighbourhood $U$ of $x$ such
that $\Gamma(U,\mathscr{B})=B$ is an algebra of finite type over
$\Gamma(U,\mathscr{O}_X)=A$. We then have $\mathscr{B}|U=\widetilde{B}$,
and for every $f\in A$, the induced $(\mathscr{O}_X|D(f))$-algebra
$\mathscr{B}|D(f)$ is of finite type, since it is isomorphic to
$\widetilde{B_f}$, and $B_f=B\otimes_A A_f$ is evidently an algebra of
finite type over $A_f$. Since the $D(f)$ form a base of the topology of
$U$, we conclude that if $\mathscr{B}$ is a quasi-coherent
$\mathscr{O}_X$-algebra of finite type, then for every open subset $V$ of
$X$, $\mathscr{B}|V$ is a quasi-coherent $(\mathscr{O}_X|V)$-algebra of
finite type.
\end{env}

\begin{proposition}[9.6.3]
\label{I.9.6.3}
Let $X$ be a locally Noetherian prescheme. Then every quasi-coherent
$\mathscr{O}_X$-algebra of finite type $\mathscr{B}$ is a coherent sheaf of
rings (\hyperref[0.5.3.7]{0, 5.3.7}).
\end{proposition}

\begin{proof}
We may again restrict to the case where $X$ is an affine scheme with
Noetherian ring $A$, and where $\mathscr{B}=\widetilde{B}$, $B$ being an
$A$-algebra of finite type; $B$ is therefore a Noetherian ring. That being
so, we must prove that the kernel $\mathscr{N}$ of a
$\mathscr{B}$-homomorphism $\mathscr{B}^m\to\mathscr{B}$ is a
$\mathscr{B}$-
\oldpage[I]{180}
module of finite type; now, it is isomorphic (as a
$\mathscr{B}$-module) to $\widetilde{N}$, where $N$ is the kernel of the
corresponding homomorphism of $B$-modules $B^m\to B$
(\hyperref[I.1.3.13]{1.3.13}). Since $B$ is Noetherian, the submodule
$N$ of $B^m$ is a $B$-module of finite type, hence there exists a
homomorphism $B^p\to B^m$ with image $N$; the sequence $B^p\to B^m\to B$
being exact, so is the corresponding sequence
$\mathscr{B}^p\to\mathscr{B}^m\to\mathscr{B}$
(\hyperref[I.1.3.5]{1.3.5}), and since $\mathscr{N}$ is the image of
$\mathscr{B}^p\to\mathscr{B}^m$
(\hyperref[I.1.3.9]{1.3.9, (i)}), the proposition is proved.
\end{proof}

\begin{corollary}[9.6.4]
\label{I.9.6.4}
Under the hypotheses of (\hyperref[I.9.6.3]{9.6.3}), for a
$\mathscr{B}$-module $\mathscr{F}$ to be coherent, it is necessary and
sufficient that it be a quasi-coherent $\mathscr{O}_X$-module and a
$\mathscr{B}$-module of finite type. If this is so, and if $\mathscr{G}$
is a sub-$\mathscr{B}$-module or a quotient $\mathscr{B}$-module of
$\mathscr{F}$, then for $\mathscr{G}$ to be a coherent $\mathscr{B}$-module,
it is necessary and sufficient that $\mathscr{G}$ be a quasi-coherent
$\mathscr{O}_X$-module.
\end{corollary}

\begin{proof}
Taking into account (\hyperref[I.9.6.1]{9.6.1}), the conditions on
$\mathscr{F}$ are evidently necessary; let us show that they are
sufficient. We may restrict to the case where $X$ is affine with Noetherian
ring $A$, $\mathscr{B}=\widetilde{B}$, where $B$ is an $A$-algebra of
finite type, $\mathscr{F}=\widetilde{M}$, where $M$ is a $B$-module, and
where there exists a surjective $\mathscr{B}$-homomorphism
$\mathscr{B}^m\to\mathscr{F}\to 0$. We then have a corresponding exact
sequence $B^m\to M\to 0$, hence $M$ is a $B$-module of finite type;
moreover, the kernel $P$ of the homomorphism $B^m\to M$ is then a
$B$-module of finite type, since $B$ is Noetherian. We conclude
(\hyperref[I.1.3.13]{1.3.13}) that $\mathscr{F}$ is the cokernel of a
$\mathscr{B}$-homomorphism $\mathscr{B}^n\to\mathscr{B}^m$, and is therefore
coherent since $\mathscr{B}$ is a coherent sheaf of rings
(\hyperref[0.5.3.4]{0, 5.3.4}). The same reasoning shows that a
quasi-coherent sub-$\mathscr{B}$-module (resp. quotient $\mathscr{B}$-module)
of $\mathscr{F}$ is of finite type, whence the second part of the
corollary.
\end{proof}

\begin{proposition}[9.6.5]
\label{I.9.6.5}
Let $X$ be a quasi-compact scheme or a prescheme whose underlying space
is Noetherian. For every quasi-coherent $\mathscr{O}_X$-algebra of finite
type $\mathscr{B}$, there exists a quasi-coherent sub-$\mathscr{O}_X$-module
of finite type $\mathscr{E}$ of $\mathscr{B}$ such that $\mathscr{E}$
generates (\hyperref[0.4.1.4]{0, 4.1.4}) the $\mathscr{O}_X$-algebra
$\mathscr{B}$.
\end{proposition}

\begin{proof}
Indeed, by hypothesis there exists a finite covering $(U_i)$ of $X$
consisting of affine open subsets such that $\Gamma(U_i,\mathscr{B})=B_i$ is
an algebra of finite type over $\Gamma(U_i,\mathscr{O}_X)=A_i$; let $E_i$ be
a sub-$A_i$-module of finite type of $B_i$ generating the $A_i$-algebra
$B_i$; by virtue of (\hyperref[I.9.4.7]{9.4.7}), there exists a
sub-$\mathscr{O}_X$-module $\mathscr{E}_i$ of $\mathscr{B}$,
quasi-coherent and of finite type, such that
$\mathscr{E}_i|U_i=\widetilde{E_i}$. It is clear that the sum
$\mathscr{E}$ of the $\mathscr{E}_i$ answers the question.
\end{proof}

\begin{proposition}[9.6.6]
\label{I.9.6.6}
Let $X$ be a prescheme whose underlying space is locally Noetherian,
or a quasi-compact scheme. Then every quasi-coherent $\mathscr{O}_X$-algebra
$\mathscr{B}$ is the inductive limit of its quasi-coherent
sub-$\mathscr{O}_X$-algebras of finite type.
\end{proposition}

\begin{proof}
Indeed, it follows from (\hyperref[I.9.4.9]{9.4.9}) that $\mathscr{B}$
is the inductive limit (as an $\mathscr{O}_X$-module) of its quasi-coherent
sub-$\mathscr{O}_X$-modules of finite type; the latter generate
quasi-coherent sub-$\mathscr{O}_X$-algebras of finite type of $\mathscr{B}$
(\hyperref[I.1.3.14]{1.3.14}), of which $\mathscr{B}$ is \emph{a
fortiori} the inductive limit.
\end{proof}
